\documentclass{exam}
\usepackage{tasks}

\usepackage{amsmath,amssymb,amsfonts}
\usepackage{graphicx}
\usepackage{amsthm}
\usepackage[french]{babel}
\usepackage{enumitem}
\newenvironment{clist}{\begin{center}\begin{inparaenum}}{\end{inparaenum}\end{center}}

\pagestyle{empty}
\graphicspath{ {../Images/} }

\usepackage{geometry}
\geometry{left=1.5cm, right=1.5cm, top=1.5cm, bottom=1.5cm}

\makeatletter
\def\gap#1{%
  \setbox\@tempboxa\hbox{{\color{white}#1}}%
  \@tempdima\wd\@tempboxa
  \rlap{\unhbox\@tempboxa}%
  \hb@xt@\@tempdima{\dotfill}%
}
\makeatother   

\theoremstyle{definition}
\newtheorem{exercice}{Exercice}

\begin{document}
\large

\noindent{\bf Contrôle 4 (B) -- 55min. \hfill Seconde générale \hfill 26/11/2024}% 

\vspace{2mm}

\noindent{\makebox[\textwidth]{Nom et prénom:\enspace\hrulefill}}

\vspace{2mm}

\begin{exercice}
Développer les expressions suivantes :
\begin{tasks}(2)
    \task $(3x - 5)^2$
    %\task $(3x - 2)(x + 4)$
    \task $3(x + 2)(2x - 3)$
\end{tasks}
\begin{center}
    \framebox[1\textwidth]{\parbox{5.5in}{%\centering \underline{Réponse:}
    \vspace{11em}}}
\end{center}
\end{exercice}

\begin{exercice}
Factoriser les expressions suivantes :
\begin{tasks}(2)
    \task $x(x + 2) + 4x$
    %\task $x(2x + 4) - x(x + 1)$
    \task $x(x + 2) + (x + 2)(x - 5)$
\end{tasks}
\begin{center}
    \framebox[1\textwidth]{\parbox{5.5in}{%\centering \underline{Réponse:}
    \vspace{11em}}}
\end{center}
\end{exercice}

\begin{exercice}
Résoudre les équations suivantes :
\begin{tasks}(3)
    \task $(2x + 6)(x - 2) = 0$
    \task $5x^2 - 20 = 0$
    \task $(2x^2 - 4)(2x - 5) = 0$
\end{tasks}
\begin{center}
    \framebox[1\textwidth]{\parbox{5.5in}{%\centering \underline{Réponse:}
    \vspace{20em}}}
\end{center}
\end{exercice}

\begin{exercice}
Écrire les expressions suivantes sous la forme $a\sqrt{b}$ :
\begin{tasks}(3)
    \task $\sqrt{192}$
    \task $\sqrt{27} + \sqrt{75} - \sqrt{12} $
    \task $\dfrac{\sqrt{162}}{\sqrt{3
    }}$
\end{tasks}
\begin{center}
    \framebox[1\textwidth]{\parbox{5.5in}{%\centering \underline{Réponse:}
    \vspace{20em}}}
\end{center}
\end{exercice}

\begin{exercice}
Écrire les expressions suivantes avec un dénominateur entier :
\begin{tasks}(2)
    \task $\dfrac{2}{\sqrt{3}}$
    \task $\dfrac{2}{\sqrt{3} + 1}$
\end{tasks}
\begin{center}
    \framebox[1\textwidth]{\parbox{5.5in}{%\centering \underline{Réponse:}
    \vspace{10em}}}
\end{center}
\end{exercice}

\end{document}
